\documentclass[tikz,border=12pt]{standalone}
\usepackage{luatexja}
\usepackage{amsmath,amssymb}
\usetikzlibrary{shapes.geometric,arrows.meta,positioning,calc,fit,backgrounds,shadows}

\begin{document}
\begin{tikzpicture}[
    font=\small,
    >=Stealth,
    node distance=1.2cm and 1.5cm,
    stepbox/.style={rectangle, rounded corners=4pt, draw=blue!70!black, fill=blue!5, thick, text width=3.4cm, align=center, inner sep=6pt, drop shadow},
    barrier/.style={rectangle, rounded corners=4pt, draw=red!70!black, fill=red!6, thick, text width=3.8cm, align=center, inner sep=6pt, drop shadow},
    solution/.style={rectangle, rounded corners=4pt, draw=teal!70!black, fill=teal!6, thick, text width=4.5cm, align=center, inner sep=6pt, drop shadow},
    arrow/.style={->, very thick, color=blue!60!black},
    alertarrow/.style={->, very thick, color=red!70!black},
    solarrow/.style={->, very thick, color=teal!70!black}
]

% Title
\node[font=\large\bfseries, text=blue!80!black] (title) at (6, 4.2) {横浜国大SPIRAL：大規模IoT診断サービス運用の教訓（USENIX Security 2025）};

% Stage 1: Detection
\node[stepbox] (scan) at (1.5, 2.5) {
    \textbf{ステップ1: 脅威検知}\\
    \footnotesize インターネット広域探索\\ポート開放・感染兆候検知\\(総務省NOTICE / 診断)
};

% Stage 2: Notification
\node[stepbox] (notify) at (6, 2.5) {
    \textbf{ステップ2: ユーザー通知}\\
    \footnotesize ISP経由の警告メール\\Web注意喚起バナー提示\\対策手順書の案内
};

% Stage 3: Remediation
\node[stepbox] (remedy) at (10.5, 2.5) {
    \textbf{ステップ3: 対策・復旧}\\
    \footnotesize パスワード変更\\ファームウェア更新\\感染マルウェアの駆除
};

% Barrier box
\node[barrier] (human) at (6, 0.2) {
    \textbf{現実の壁（ヒューマンファクター）}\\
    \footnotesize ・警告を不審メールと誤認\\・管理画面の開き方が不明\\・専門用語（IP/ポート）の難解さ\\・機器EOLでパッチ未提供
};

% Success vs Dropout
\node[barrier] (dropout) at (10.5, 0.2) {
    \textbf{高い離脱率 (Dropout)}\\
    \footnotesize 実際に設定変更・駆除に\\至るユーザーはごく一部\\（自己責任モデルの限界）
};

% Future solution
\node[solution] (future) at (6, -2.5) {
    \textbf{学術的知見から導かれた解決策（佐々木・吉岡研）}\\
    \footnotesize ① \textbf{Secure-by-Design}: 初期PW廃止・初回強制設定\\
    ② \textbf{自動更新}: ユーザー無操作での自動セキュリティパッチ\\
    ③ \textbf{ネットワーク協調}: ISP側での悪性通信フィルタリング
};

% Arrows
\draw[arrow] (scan) -- node[above, font=\scriptsize] {診断結果連携} (notify);
\draw[arrow] (notify) -- node[above, font=\scriptsize] {対策要請} (remedy);
\draw[alertarrow] (notify) -- (human);
\draw[alertarrow] (human) -- (dropout);
\draw[alertarrow] (remedy) -- (dropout);
\draw[solarrow] (human) -- node[right, font=\scriptsize] {構造改革} (future);

\end{tikzpicture}
\end{document}
